\chapter{Three-dimensional rotations and orientation representations}
\label{chap:lecture04}

A robot that moves outside a plane requires a spatial description of
orientation. One angle no longer suffices: a body can turn about three
independent directions. The construction of a rotation matrix remains the
same, however. Its columns are the unit axes of one frame expressed in another.
In three dimensions, these columns have three components, giving a $3\times3$
matrix.

This chapter constructs rotations about individual axes, combines them into
a yaw--pitch--roll sequence, recovers those angles from a rotation matrix,
and introduces unit quaternions as a practical orientation representation.
All frames used for point-coordinate transformations here share an origin;
translation is set aside so that the orientation conventions can be developed
explicitly.

\section{Spatial frames and rotation matrices}
Let $W$ be a right-handed world frame and $R$ a right-handed robot frame.
Each has three mutually perpendicular unit axes. Write the robot's axes in
world components as $[\boldsymbol{x}_R]^W$, $[\boldsymbol{y}_R]^W$, and
$[\boldsymbol{z}_R]^W$. Then
\begin{equation}
 R_R^W=\begin{bmatrix}
 [\boldsymbol{x}_R]^W &[\boldsymbol{y}_R]^W &[\boldsymbol{z}_R]^W
 \end{bmatrix}\in\R^{3\times3}.
 \label{eq:l04-columns}
\end{equation}
As in the planar case, the subscript gives the input frame and the superscript
gives the output frame of the coordinate map. For a point $B$ and common origins,
\begin{equation}
 B^W=R_R^W B^R,\qquad
 B^R=\begin{bmatrix}x_B^R\\y_B^R\\z_B^R\end{bmatrix}.
\end{equation}
The same rotation maps the components of a direction vector between frames
regardless of the locations of their origins.

Positive elementary angles follow the right-hand rule: point the right thumb
along the positive rotation axis and curl the fingers in the positive direction.
Viewed from the positive end of that axis toward the origin, the positive
rotation appears counterclockwise. A perspective drawing may distort lengths
and visible angles; it does not change this convention.

\section{Elementary rotations}
An elementary rotation leaves one axis fixed and rotates the perpendicular
plane. The three basic cases are shown in Figure~\ref{fig:l04-elementary}.
Each is constructed by expressing the rotated basis vectors in the original
frame, then placing those vectors in columns.

\begin{figure}[htbp]
\centering% Three exact 3D constructions projected onto the page; lengths in the
% projection are not physical lengths. The red axes are matrix columns.
\begin{tikzpicture}[x={(-.65cm,-.4cm)},y={(.85cm,-.25cm)},z={(0cm,.95cm)}]
\foreach \panel/\name in {0/z,4.8/y,9.6/x} {
\begin{scope}[shift={(\panel cm,0cm)}]
\draw[axis] (0,0,0)--(1.8,0,0) node[below left] {$x_W$};
\draw[axis] (0,0,0)--(0,1.8,0) node[right] {$y_W$};
\draw[axis] (0,0,0)--(0,0,1.8) node[above] {$z_W$};
\if\name z
 \draw[axis,robotframe] (0,0,0)--({1.5*cos(35)},{1.5*sin(35)},0) node[below] {$x_R$};
 \draw[axis,robotframe] (0,0,0)--({-1.5*sin(35)},{1.5*cos(35)},0) node[right] {$y_R$};
 \draw[axis,robotframe] (0,0,0)--(0,0,1.5) node[left] {$z_R$};
\fi
\if\name y
 \draw[axis,robotframe] (0,0,0)--({1.5*cos(35)},0,{-1.5*sin(35)}) node[below] {$x_R$};
 \draw[axis,robotframe] (0,0,0)--(0,1.5,0) node[above] {$y_R$};
 \draw[axis,robotframe] (0,0,0)--({1.5*sin(35)},0,{1.5*cos(35)}) node[left] {$z_R$};
\fi
\if\name x
 \draw[axis,robotframe] (0,0,0)--(1.5,0,0) node[above left] {$x_R$};
 \draw[axis,robotframe] (0,0,0)--(0,{1.5*cos(35)},{1.5*sin(35)}) node[right] {$y_R$};
 \draw[axis,robotframe] (0,0,0)--(0,{-1.5*sin(35)},{1.5*cos(35)}) node[left] {$z_R$};
\fi
\node at (0,-1.7,0) {};
\node[anchor=north] at (0cm,-1.8cm) {$R_{\name}$: rotation about $\name$};
\end{scope}
}
\end{tikzpicture}

\caption{Elementary rotations in perspective. Black axes are the world frame;
red axes are the rotated frame. The rotation axis is shared by the two frames.
The three panels are separate constructions, each with a positive angle.}
\label{fig:l04-elementary}
\end{figure}

\subsection{Rotation about the z axis}
Let $\psi$ denote a rotation about $z_W$. The $z$ axis remains unchanged, while
the $x$ and $y$ axes rotate in the $xy$ plane:
\[
 [\boldsymbol{x}_R]^W=\begin{bmatrix}\cos\psi\\\sin\psi\\0\end{bmatrix},\quad
 [\boldsymbol{y}_R]^W=\begin{bmatrix}-\sin\psi\\\cos\psi\\0\end{bmatrix},\quad
 [\boldsymbol{z}_R]^W=\begin{bmatrix}0\\0\\1\end{bmatrix}.
\]
Therefore,
\begin{equation}
 \boxed{R_z(\psi)=\begin{bmatrix}
 \cos\psi&-\sin\psi&0\\\sin\psi&\cos\psi&0\\0&0&1
 \end{bmatrix}.}
 \label{eq:l04-rz}
\end{equation}
The upper-left block is exactly the planar rotation matrix. The last coordinate
is unchanged because the rotation takes place in the plane perpendicular to $z$.

\subsection{Rotation about the y axis}
For a positive angle $\theta$ about $y_W$, the $x$ axis turns toward negative
$z$ and the $z$ axis turns toward positive $x$. The basis columns are
\[
 [\boldsymbol{x}_R]^W=\begin{bmatrix}\cos\theta\\0\\-\sin\theta\end{bmatrix},\quad
 [\boldsymbol{y}_R]^W=\begin{bmatrix}0\\1\\0\end{bmatrix},\quad
 [\boldsymbol{z}_R]^W=\begin{bmatrix}\sin\theta\\0\\\cos\theta\end{bmatrix}.
\]
Thus
\begin{equation}
 \boxed{R_y(\theta)=\begin{bmatrix}
 \cos\theta&0&\sin\theta\\0&1&0\\-\sin\theta&0&\cos\theta
 \end{bmatrix}.}
 \label{eq:l04-ry}
\end{equation}
The signs follow from the right-hand rule. The two-dimensional pattern is
visible in the ordered plane $(z,x)$, where positive rotation carries $z$
toward $x$; using the order $(x,z)$ without accounting for orientation would
reverse the signs.

\subsection{Rotation about the x axis}
For a positive angle $\varphi$ about $x_W$, the $x$ axis is fixed and the
$y$ axis turns toward positive $z$. Hence
\[
 [\boldsymbol{x}_R]^W=\begin{bmatrix}1\\0\\0\end{bmatrix},\quad
 [\boldsymbol{y}_R]^W=\begin{bmatrix}0\\\cos\varphi\\\sin\varphi\end{bmatrix},\quad
 [\boldsymbol{z}_R]^W=\begin{bmatrix}0\\-\sin\varphi\\\cos\varphi\end{bmatrix},
\]
and
\begin{equation}
 \boxed{R_x(\varphi)=\begin{bmatrix}
 1&0&0\\0&\cos\varphi&-\sin\varphi\\0&\sin\varphi&\cos\varphi
 \end{bmatrix}.}
 \label{eq:l04-rx}
\end{equation}

\begin{figure}[htbp]
\centering\begin{tikzpicture}[scale=.8]
\foreach \shift/\horizontal/\vertical/\angle/\moving in {
0/x_W/y_W/\psi/x_R,5.1/z_W/x_W/\theta/z_R,10.2/y_W/z_W/\varphi/y_R} {
\begin{scope}[shift={(\shift,0)}]
\draw[axis] (0,0)--(2.7,0) node[right] {$\horizontal$};
\draw[axis] (0,0)--(0,2.5) node[above] {$\vertical$};
\draw[axis,robotframe] (0,0)--(35:2.2) node[above right] {$\moving$};
\draw[projection] (35:2.2)--({2.2*cos(35)},0);
\draw[projection] (35:2.2)--(0,{2.2*sin(35)});
\node[below] at (.9,0) {$\cos\angle$};
\node[left] at (0,.65) {$\sin\angle$};
\draw[-{Stealth}] (.65,0) arc (0:35:.65);
\node at (17.5:.95) {$\angle$};
\end{scope}
}
\end{tikzpicture}

\caption{Plane views of the elementary rotations. From left to right, the
ordered planes are $(x_W,y_W)$, $(z_W,x_W)$, and $(y_W,z_W)$. Each red
vector has unit length; its projections explain the sine and cosine entries.}
\label{fig:l04-planes}
\end{figure}

A quick sign check is to use a quarter-turn: $R_z(\pi/2)$ maps the first
basis vector to the second, $R_y(\pi/2)$ maps the third to the first, and
$R_x(\pi/2)$ maps the second to the third. At zero angle, all three matrices
reduce to $I_3$.

\section{What makes a matrix a valid spatial rotation?}
Write $R=[u\;v\;w]$. Its columns must be unit vectors and mutually
perpendicular:
\begin{equation}
 u\trans u=v\trans v=w\trans w=1,\qquad
 u\trans v=u\trans w=v\trans w=0.
 \label{eq:l04-constraints}
\end{equation}
Together these six scalar constraints give $R\trans R=I_3$. They constrain
the nine matrix entries to three continuous degrees of freedom. To describe
a rotation between right-handed frames, rather than a reflection, one also
requires $\det R=+1$. The determinant condition selects the proper-rotation
component; it does not remove another continuous degree of freedom.

For a valid rotation,
\begin{equation}
 R^{-1}=R\trans,\qquad R_F^A=(R_A^F)\trans.
\end{equation}
Composition still follows the frame chain, $R_C^A=R_B^A R_C^B$, and lengths
are preserved. These are the same structural properties used in two dimensions.

\subsection{An illustrative set of components}
To illustrate how columns are assembled, consider the rough component array
\begin{equation}
 M=\begin{bmatrix}.8&-.3&.4\\.2&.7&.15\\-.1&.2&.7\end{bmatrix}.
 \label{eq:l04-rough}
\end{equation}
Its first column describes a direction with a large positive $x$ component,
a smaller positive $y$ component, and a negative $z$ component. The other
columns can be read in the same way. However, these illustrative numbers do
\emph{not} define a rotation matrix. For example,
\[
 u\trans u=.8^2+.2^2+(-.1)^2=.69\ne1,\qquad
 u\trans v=.8(-.3)+.2(.7)+(-.1)(.2)=-.12\ne0.
\]
Even normalizing each column separately would not make the columns mutually
perpendicular. A valid frame must satisfy all the orthonormality conditions
and the right-handedness condition, not just have plausible-looking components.

\section{Roll, pitch, and yaw as a sequence of rotations}
A general spatial orientation can be represented by a sequence of three
elementary rotations. Many sequences are possible, and the convention must
be specified before interpreting three angle values. Here the convention is
the \emph{intrinsic Z--Y--X sequence}: each successive rotation is about an
axis of the frame obtained after the previous rotation.

Starting from frame $W$, define intermediate frames 1 and 2 and final frame 3:
\begin{enumerate}
 \item Rotate by yaw $\psi$ about $z_W$ to obtain frame 1.
 \item Rotate by pitch $\theta$ about the new axis $y_1$ to obtain frame 2.
 \item Rotate by roll $\varphi$ about the new axis $x_2$ to obtain frame 3.
\end{enumerate}
Although the angles are often listed as roll, pitch, yaw, the construction
above proceeds through yaw, pitch, then roll.

\begin{figure}[htbp]
\centering\begin{tikzpicture}
\foreach \x/\frame in {0/W,3/1,6/2,9/3} {
 \node[draw,rounded corners,minimum width=1cm,minimum height=.8cm] (f\frame) at (\x,0) {Frame $\frame$};
}
\draw[axis] (fW)--(f1) node[midway,above] {yaw $\psi$} node[midway,below] {about $z_W$};
\draw[axis] (f1)--(f2) node[midway,above] {pitch $\theta$} node[midway,below] {about $y_1$};
\draw[axis] (f2)--(f3) node[midway,above] {roll $\varphi$} node[midway,below] {about $x_2$};
\node at (1.5,-1.15) {$z_1=z_W$};
\node at (4.5,-1.15) {$y_2=y_1$};
\node at (7.5,-1.15) {$x_3=x_2$};
\end{tikzpicture}

\caption{The intrinsic Z--Y--X construction. The axis used at each stage stays
fixed during that stage. Arrows show the sequence of frame construction;
coordinate maps from frame 3 back to $W$ act in the reverse traversal.}
\label{fig:l04-sequence}
\end{figure}

The elementary frame relationships are
\[
 R_1^W=R_z(\psi),\qquad R_2^1=R_y(\theta),\qquad R_3^2=R_x(\varphi).
\]
Applying the frame-chain rule gives
\begin{equation}
 \boxed{R_3^W=R_1^W R_2^1 R_3^2
 =R_z(\psi)R_y(\theta)R_x(\varphi).}
 \label{eq:l04-rpy}
\end{equation}
This product maps coordinates from frame 3 to $W$. When it multiplies $B^3$,
the rightmost factor acts first: $B^3\to B^2\to B^1\to B^W$. This algebraic
order is consistent with constructing the frames successively about their
moving axes; it should not be confused with three rotations about fixed world
axes in the same chronological order.

\subsection{Order matters in three dimensions}
Unlike planar rotations, rotations about different spatial axes do not
generally commute. With $e_x=[1\;0\;0]\trans$,
\[
 R_z(\pi/2)R_x(\pi/2)e_x=e_y,\qquad
 R_x(\pi/2)R_z(\pi/2)e_x=e_z.
\]
The first product leaves $e_x$ unchanged in its first step, then turns it
toward $e_y$. The second turns it toward $e_y$ first, then toward $e_z$.
Thus changing the factor order generally changes the orientation.

\subsection{Expanded form}
For compactness, write $c_a=\cos a$ and $s_a=\sin a$. Multiplying the three
elementary matrices yields
\begin{equation}
 R=\begin{bmatrix}
 c_\psi c_\theta & c_\psi s_\theta s_\varphi-s_\psi c_\varphi & c_\psi s_\theta c_\varphi+s_\psi s_\varphi\\
 s_\psi c_\theta & s_\psi s_\theta s_\varphi+c_\psi c_\varphi & s_\psi s_\theta c_\varphi-c_\psi s_\varphi\\
 -s_\theta & c_\theta s_\varphi & c_\theta c_\varphi
 \end{bmatrix}.
 \label{eq:l04-expanded}
\end{equation}
The factored form is often easier to evaluate and check. The expanded form is
particularly useful for identifying which entries contain the information
needed to recover the angles.

\section{Worked example}
Let the roll, pitch, and yaw angles be
\[
 \varphi=\frac{\pi}{4},\qquad \theta=\frac{\pi}{3},\qquad \psi=\frac{\pi}{2}.
\]
Then
\begin{align}
 R_3^W
 &=\begin{bmatrix}0&-1&0\\1&0&0\\0&0&1\end{bmatrix}
 \begin{bmatrix}1/2&0&\sqrt3/2\\0&1&0\\-\sqrt3/2&0&1/2\end{bmatrix}
 \begin{bmatrix}1&0&0\\0&\sqrt2/2&-\sqrt2/2\\0&\sqrt2/2&\sqrt2/2\end{bmatrix}\\
 &=\begin{bmatrix}
 0&-\sqrt2/2&\sqrt2/2\\
 1/2&\sqrt6/4&\sqrt6/4\\
 -\sqrt3/2&\sqrt2/4&\sqrt2/4
 \end{bmatrix}.
 \label{eq:l04-example}
\end{align}
For instance, the first column has squared length $1/4+3/4=1$ and has zero
dot product with the second column. More generally, the product of these
three proper rotation matrices is itself a proper rotation matrix.

If $B^3=[1\;0\;0]\trans$, its world coordinates are the first column of
$R_3^W$. To return to frame 3, use $(R_3^W)\trans$. Thus the orientation
representation serves directly as a coordinate transformation.

\section{Recovering roll, pitch, and yaw}
Let a valid rotation matrix be written as
\[
 R=\begin{bmatrix}r_{11}&r_{12}&r_{13}\\r_{21}&r_{22}&r_{23}\\r_{31}&r_{32}&r_{33}\end{bmatrix}.
\]
For the convention in Equation~\eqref{eq:l04-rpy}, choose the principal pitch
range $-\pi/2<\theta<\pi/2$. In this nonsingular range,
\begin{equation}
 \boxed{\begin{aligned}
 \varphi&=\operatorname{atan2}(r_{32},r_{33}),\\
 \theta&=\operatorname{atan2}\!\left(-r_{31},\sqrt{r_{32}^2+r_{33}^2}\right),\\
 \psi&=\operatorname{atan2}(r_{21},r_{11}).
 \end{aligned}}
 \label{eq:l04-extraction}
\end{equation}
The entries in Equation~\eqref{eq:l04-expanded} explain these formulas:
\[
 (r_{32},r_{33})=c_\theta(s_\varphi,c_\varphi),\qquad
 (r_{21},r_{11})=c_\theta(s_\psi,c_\psi).
\]
Since $c_\theta>0$ on the chosen range, each pair has the same direction as
the corresponding sine--cosine pair. Also,
\[
 -r_{31}=s_\theta,\qquad
 \sqrt{r_{32}^2+r_{33}^2}=|c_\theta|=c_\theta.
\]
Applied to Equation~\eqref{eq:l04-example}, the formulas give
\[
 \varphi=\operatorname{atan2}(\sqrt2/4,\sqrt2/4)=\pi/4,
 \quad \theta=\operatorname{atan2}(\sqrt3/2,1/2)=\pi/3,
\]
and $\psi=\operatorname{atan2}(1/2,0)=\pi/2$, recovering the original angles.

\subsection{Why atan2 takes two arguments}
The function $\operatorname{atan2}(y,x)$ returns the angle of the pair $(x,y)$,
using both signs to identify the quadrant. The argument order here is
\emph{vertical component first, horizontal component second}. The numeral 2
is part of the function's name, not a multiplication or exponent.

The ratio $y/x$ alone loses information: $(1,1)$ and $(-1,-1)$ both have ratio
one, but point in opposite directions. Ordinary $\arctan(y/x)$ returns the
same principal angle for both, whereas
\[
 \operatorname{atan2}(1,1)=\pi/4,\qquad
 \operatorname{atan2}(-1,-1)=-3\pi/4.
\]
Unlike division followed by arctangent, $\operatorname{atan2}$ also handles
$x=0$ with $y\ne0$, giving $\pi/2$ or $-\pi/2$.

\begin{figure}[htbp]
\centering\begin{tikzpicture}[scale=1.65]
\draw[axis] (-1.5,0)--(1.65,0) node[right] {$x$};
\draw[axis] (0,-1.4)--(0,1.5) node[above] {$y$};
\draw[gray] (0,0) circle (1);
\draw[axis,sensorframe] (0,0)--(45:1.25) node[above right] {$(1,1)$};
\draw[axis,robotframe] (0,0)--(225:1.25) node[below left] {$(-1,-1)$};
\draw[-{Stealth},sensorframe] (.55,0) arc (0:45:.55);
\node[sensorframe] at (25:.85) {$\pi/4$};
\draw[-{Stealth},robotframe] (.38,0) arc (0:-135:.38);
\node[robotframe] at (-65:.85) {$-3\pi/4$};
\end{tikzpicture}

\caption{Equal ratios do not imply equal directions. The signs of both
components distinguish the first and third quadrants. The circle indicates
direction only; the two labeled vectors need not have unit length.}
\label{fig:l04-atan2}
\end{figure}

Tangent has period $\pi$, while a full revolution has length $2\pi$.
The principal arctangent takes values in $(-\pi/2,\pi/2)$; a conventional
principal range for $\operatorname{atan2}$ spans $2\pi$, such as $(-\pi,\pi]$.
At the negative horizontal axis, $\pi$ and $-\pi$ represent the same direction;
endpoint handling can vary between implementations. At $(x,y)=(0,0)$ there is
no geometric direction, whatever numerical value a software routine returns.

\subsection{Singular orientations and angle ranges}
The three-angle representation covers all orientations, but is not globally
unique or nonsingular. At $\theta=\pm\pi/2$, the pairs used to extract yaw
and roll both vanish:
\[
 r_{32}=r_{33}=r_{21}=r_{11}=0.
\]
Consequently, Equation~\eqref{eq:l04-extraction} cannot determine yaw and roll
separately. At $\theta=\pi/2$, for example, the matrix depends on their
difference $\varphi-\psi$. This is the Euler-angle singularity often called
\emph{gimbal lock}. It is a limitation of the angle coordinates, not a loss
of information in the rotation matrix.

An implementation must choose a convention for the singular case rather than
treat $\operatorname{atan2}(0,0)$ as a uniquely meaningful angle. Near this
case, small matrix perturbations can produce large changes in extracted yaw
and roll. Away from it, the principal ranges are typically roll and yaw in
a $2\pi$ interval, and pitch in $[-\pi/2,\pi/2]$.

\section{Choosing and maintaining an orientation representation}
Angles make individual turns easy to describe, while matrices make coordinate
transformation and composition direct. The map $B^W=R_R^W B^R$ is linear in
the coordinates of $B$ for a fixed rotation; the dependence of $R_R^W$ on the
angles remains trigonometric.

Finite-precision arithmetic can cause repeatedly multiplied rotation matrices
to drift from orthogonality. Useful checks are whether $R\trans R$ remains
close to $I_3$ and whether $\det R$ remains close to $+1$. Restoring unit
column lengths alone is insufficient: perpendicularity and handedness must
also be retained. A correction procedure must restore a proper orthonormal frame.

Angle representations have their own difficulties, including wrapping at the
ends of the chosen intervals and singular orientations. Unit quaternions
provide another representation used in robotics software. No representation
removes the need to specify conventions and manage numerical error.

\subsection{Unit quaternions in practice}
A unit quaternion represents orientation using four real components. With
the scalar component first, write
\begin{equation}
 q=\begin{bmatrix}q_w&q_x&q_y&q_z\end{bmatrix}\trans,
 \qquad q\trans q=q_w^2+q_x^2+q_y^2+q_z^2=1.
 \label{eq:l04-quaternion}
\end{equation}
Four components with one unit-norm constraint leave three continuous degrees
of freedom, just as for a rotation matrix. The quaternions $q$ and $-q$
represent the same orientation; this discrete ambiguity does not add a degree
of freedom. Component order must always be checked when exchanging data:
some software places the scalar component last.

For a nonzero quaternion affected by small numerical drift, restoring unit
norm is straightforward:
\[
 q\leftarrow\frac{q}{\|q\|},\qquad
 \|q\|=\sqrt{q_w^2+q_x^2+q_y^2+q_z^2}.
\]
A zero quaternion cannot be normalized and does not represent an orientation.
Normalization enforces the representation's constraint; it does not undo
every numerical error in the represented orientation.

The four components are less immediately intuitive than roll, pitch, and yaw,
but are convenient for storing and exchanging orientations. A useful workflow
is to retain a unit quaternion in software, convert to a rotation matrix for
coordinate calculations, and extract angles when interpreting the result.
To compose two orientations using familiar matrix algebra, convert both to
matrices, multiply in the appropriate frame order, and convert back if needed.
Direct quaternion composition is also possible and efficient; its algebra is
outside the scope of this introduction.

Appendix~\ref{app:rotation-conversions} gives all six conversion directions
between matrices, RPY angles, and scalar-first unit quaternions, including
their validity conditions. For example, a quaternion can be converted to
RPY angles to inspect yaw, while remembering that the chosen angle
representation still has its pitch singularities.

\section{Key identities}
\begin{center}
\renewcommand{\arraystretch}{1.3}
\begin{tabular}{ll}
\toprule
Property or operation & Rule \\
\midrule
Proper rotation & $R\trans R=I_3$, $\det R=1$ \\
Coordinate transformation, common origin & $B^W=R_R^W B^R$ \\
Reverse transformation & $R_W^R=(R_R^W)\trans$ \\
Frame composition & $R_C^A=R_B^A R_C^B$ \\
Intrinsic Z--Y--X convention & $R=R_z(\psi)R_y(\theta)R_x(\varphi)$ \\
Nonsingular principal pitch range & $-\pi/2<\theta<\pi/2$ \\
\bottomrule
\end{tabular}
\end{center}

\clearpage
\section{Additional exercises}
Use right-handed frames, column vectors, and the intrinsic Z--Y--X convention.
Show your calculations and keep exact trigonometric values where possible.
\begin{enumerate}
\item \textbf{Elementary rotations.}
Write $R_z(\pi/2)$, $R_y(\pi/2)$, and $R_x(\pi/2)$. For each matrix,
identify its unchanged axis and determine the images of all three standard
basis vectors. Check the signs using the right-hand rule.

\item \textbf{Checking a candidate matrix.}
For the matrix $M$ in Equation~\eqref{eq:l04-rough}, compute the squared
lengths of its three columns and all three pairwise dot products. Explain why
normalizing the columns individually cannot produce a valid rotation matrix.

\item \textbf{A valid spatial orientation.}
Consider
\[
 R=\begin{bmatrix}0&-1&0\\1&0&0\\0&0&1\end{bmatrix},\qquad
 B^R=\begin{bmatrix}1\\2\\3\end{bmatrix}.
\]
Compute $B^W=RB^R$, recover $B^R$ using the transpose, and verify length
preservation. State which elementary rotation $R$ represents.

\item \textbf{Noncommutativity.}
Compute $R_x(\pi/2)R_y(\pi/2)$ and $R_y(\pi/2)R_x(\pi/2)$.
Apply both to $e_z=[0\;0\;1]\trans$. Explain why the two results differ.

\item \textbf{From angles to a matrix and back.}
Take $\varphi=\pi/4$, $\theta=0$, and $\psi=\pi/2$.
Compute $R=R_z(\psi)R_y(\theta)R_x(\varphi)$, then recover all three angles
with Equation~\eqref{eq:l04-extraction}. Identify the frame axis used at each
stage of the intrinsic construction.

\item \textbf{Quadrants and atan2.}
Evaluate $\operatorname{atan2}(1,-1)$,
$\operatorname{atan2}(-1,-1)$, and $\operatorname{atan2}(-1,0)$.
Sketch the corresponding directions. Explain what goes wrong if each result
is computed using only $\arctan(y/x)$.

\item \textbf{A singular orientation.}
Set $\theta=\pi/2$ in Equation~\eqref{eq:l04-expanded}. Show that the
result depends on $\varphi-\psi$. Verify that the triples
$(\varphi,\theta,\psi)=(0,\pi/2,0)$ and $(\pi/4,\pi/2,\pi/4)$ give
the same matrix. Why can the extraction formulas not recover the two triples
as distinct orientations?
\end{enumerate}
